L'algorithmique tient une place à part : elle n'est pas une branche des mathématiques mais un
endroit où elles s'appliquent à elles-mêmes. Un programme est un objet dont on peut démontrer
quelque chose — qu'il calcule ce qu'on croit, qu'il s'arrête, que deux versions donnent le
même résultat.

Ce sont ces énoncés-là que le chapitre retient, et eux seuls. « Écrire un programme qui… »
est une tâche, pas une proposition ; ce qui se démontre, c'est l'invariant d'une boucle ou
l'égalité de deux programmes de calcul pour toute entrée. La distinction est la même
qu'entre tracer une médiatrice et démontrer sa caractérisation, et elle traverse tout ce
livre.

Un assistant de preuve est, de ce point de vue, le lieu naturel du sujet : les démonstrations
de ce livre sont elles-mêmes des programmes, vérifiés par un autre programme.
